% Repository-ready section file.
% Assumes: amsmath, amssymb, amsthm, mathtools, enumitem.
% If theorem environments are not already defined upstream, uncomment the block below.
%
% \newtheorem{theorem}{Theorem}[section]
% \newtheorem{proposition}[theorem]{Proposition}
% \newtheorem{lemma}[theorem]{Lemma}
% \newtheorem{corollary}[theorem]{Corollary}
% \newtheorem{conjecture}[theorem]{Conjecture}
% \theoremstyle{definition}
% \newtheorem{definition}[theorem]{Definition}
% \newtheorem{remark}[theorem]{Remark}
% \newtheorem{problem}[theorem]{Problem}
%
% The status conventions in this file are:
%   - Theorems/Propositions/Lemmas with proofs are formal or representation-theoretic statements
%     proved in this section.
%   - Conjectures are physically motivated or modular-completion statements.
%   - The final ``provisional'' subsection records exact low-mode coefficient identities discussed in
%     the working thread, but labels them explicitly as unverified and not used elsewhere.

\providecommand{\A}{\mathcal A}
\providecommand{\B}{\mathcal B}
\providecommand{\C}{\mathbb C}
\providecommand{\D}{\mathcal D}
\providecommand{\E}{\mathcal E}
\providecommand{\F}{\mathcal F}
\providecommand{\G}{\mathfrak G}
\providecommand{\M}{\mathcal M}
\providecommand{\N}{\mathcal N}
\providecommand{\Q}{\mathbb Q}
\providecommand{\R}{\mathbb R}
\providecommand{\Z}{\mathbb Z}
\providecommand{\mc}{\mathfrak m}
\providecommand{\h}{\hbar}
\providecommand{\id}{\mathrm{id}}
\providecommand{\Hom}{\operatorname{Hom}}
\providecommand{\MC}{\operatorname{MC}}
\providecommand{\Cone}{\operatorname{Cone}}
\providecommand{\Obs}{\operatorname{Obs}}
\providecommand{\Sym}{\operatorname{Sym}}
\providecommand{\Tr}{\operatorname{Tr}}
\providecommand{\Span}{\operatorname{Span}}
\providecommand{\End}{\operatorname{End}}
\providecommand{\Ext}{\operatorname{Ext}}
\providecommand{\Defcyc}{\operatorname{Def}_{\mathrm{cyc}}}
\providecommand{\Ran}{\operatorname{Ran}}
\providecommand{\gr}{\operatorname{gr}}
\providecommand{\Hb}{H_b^{0,\bullet}(S^3)}
\providecommand{\HCR}{H_{\mathrm{CR}}^{0,\bullet}(S^3)}
\providecommand{\delbar}{\bar\partial}
\providecommand{\delbarb}{\bar\partial_b}
\providecommand{\ot}{\widehat\otimes}
\providecommand{\ov}[1]{\overline{#1}}

\section{Modular twisted/celestial holography from first principles}
% label removed: sec:modular-twisted-celestial-first-principles

\subsection{The problem, restated at its true level}
% label removed: subsec:problem-restated-true-level

The disk-level story is already visible in the existing literature.  Zeng shows that the
Kaluza--Klein reduction of a class of $6$-dimensional holomorphic theories produces a
$3$-dimensional holomorphic--topological theory whose effective interactions are obtained by
homotopy transfer, and that with a suitable transverse boundary condition the resulting
boundary chiral algebra coincides with the universal defect chiral algebra of the unreduced
model \cite{ZengTwistedHolography}.  Costello--Dimofte--Gaiotto identify the basic
bulk--boundary package in the holomorphic twist of $3$d $\mathcal N=2$ theories: the bulk
chiral algebra is commutative and carries a shifted Poisson bracket, while boundary chiral
algebras are modules for the bulk and are expected to carry higher $A_\infty$-type operations
\cite{CostelloDimofteGaiottoBoundary}.  Gaiotto--Kulp--Wu then supply the perturbative
higher-operation language in which regularized Feynman diagrams satisfy quadratic/Wess--
Zumino identities and control holomorphic--topological OPE data, including defect and
boundary operations \cite{GaiottoKulpWuHigher}.  Finally, Gui--Li--Zeng formulate chiral
quadratic duality by Maurer--Cartan theory, showing in the quadratic setting that morphisms
$A\to B$ are controlled by Maurer--Cartan elements in $A^!\otimes B$, and making precise
how curvature enters via twisted pairs and curved chiral DG structures
\cite{GuiLiZengQuadraticDuality}.

Taken together, these results do \emph{not} merely suggest a new slogan.  They force a
reorganization of the problem.

The correct object is not the boundary vertex algebra by itself.  It is the full modular
Maurer--Cartan lift of the boundary algebra, of which the familiar disk-level boundary OPE is
only the genus-$0$ truncation.  In other words:
\begin{quote}
\emph{twisted/celestial holography is naturally a modular lifting problem for the boundary
chiral algebra.}
\end{quote}

The rest of this section turns that principle into a concrete mathematical package.

\subsection{The abstract modular package}
% label removed: subsec:abstract-modular-package

We isolate the minimum formal structure needed to speak about modular completion.

\begin{definition}[Complete genus-filtered cyclic deformation package]
% label removed: def:complete-genus-filtered-package
A \emph{complete genus-filtered cyclic deformation package} consists of the following data.
\begin{enumerate}[label=\textup{(\roman*)}]
\item A complete $\h$-adically filtered graded vector space
\[
\mathfrak g^{\mathrm{mod}}=\prod_{g\ge 0}\h^g\mathfrak g^{(g)}.
\]
\item A family of multilinear operations
\[
\ell_k^{(g)}: \Sym^k(\mathfrak g^{\mathrm{mod}})\longrightarrow
\mathfrak g^{\mathrm{mod}}[2-k],
\qquad g\ge 0,\ k\ge 1,
\]
which together define a complete cyclic $L_\infty$-algebra.
\item A distinguished genus-$0$ Maurer--Cartan element
\[
\pi\in \mathfrak g^{(0),1}
\]
solving the genus-$0$ Maurer--Cartan equation
\[
\sum_{k\ge 1}\frac1{k!}\,\ell_k^{(0)}(\pi^{\otimes k})=0.
\]
\end{enumerate}
A \emph{modular lift} of $\pi$ is a formal series
\[
\Pi(\h)=\pi+\sum_{g\ge 1}\h^g\Theta^{(g)}
\]
satisfying the full Maurer--Cartan equation
\begin{equation}
% label removed: eq:full-modular-mc
\sum_{r\ge 0}\h^r\sum_{k\ge 1}\frac1{k!}\,
\ell_k^{(r)}\!\bigl(\Pi(\h)^{\otimes k}\bigr)=0.
\end{equation}
\end{definition}

\begin{remark}
In the applications we have in mind, $\mathfrak g^{\mathrm{mod}}$ is the completed cyclic
boundary deformation complex tensored with a genus package such as
$R\Gamma(\ov{\M}_{\bullet,\bullet},\Q)$, or equivalently the graph/Feynman transform of the
genus-$0$ deformation theory.  The point of Definition
\ref{def:complete-genus-filtered-package} is to separate the formal deformation theory from the
physical realization.
\end{remark}

\begin{conjecture}[Modular twisted/celestial holography; \ClaimStatusConjectured]
% label removed: conj:modular-twisted-celestial
Let $(T,\mathfrak b)$ be a twisted or celestial holographic model whose disk-level boundary
algebra $\A_\partial=\A_\partial(T,\mathfrak b)$ is identified, for a transverse boundary
condition, with the universal defect chiral algebra of the unreduced theory.  Then there exists
a complete genus-filtered cyclic deformation package
\[
\mathfrak g^{\mathrm{mod}}_\partial=\prod_{g\ge 0}\h^g\mathfrak g^{(g)}_\partial
\]
and a distinguished genus-$0$ Maurer--Cartan element
\[
\pi_\partial\in \mathfrak g_\partial^{(0),1}
\]
such that:
\begin{enumerate}[label=\textup{(\roman*)}]
\item the disk-level boundary OPE and transferred $A_\infty$-module structure are encoded by
$\pi_\partial$;
\item the full twisted/celestial holographic dictionary is equivalent to a modular lift
\[
\Pi_\partial(\h)=\pi_\partial+\sum_{g\ge 1}\h^g\Theta_\partial^{(g)}
\]
of $\pi_\partial$;
\item the coefficients $\Theta_\partial^{(g)}$ are the genus-$g$/loop/$1/N$ corrections to the
bulk--boundary dictionary;
\item under the appropriate bulk--boundary Koszul/Verdier comparison, this modular lift is the
same deformation problem as the corresponding bulk modular lift.
\end{enumerate}
Equivalently: the known boundary chiral algebra is the genus-$0$ shadow of a canonical
modular Maurer--Cartan element.
\end{conjecture}

\subsection{Order-by-order obstruction theory}
% label removed: subsec:order-by-order-obstruction-theory

The first substantial gain from the modular viewpoint is that higher-genus corrections stop
being ad hoc and become obstruction classes in a precise cohomology.

\begin{definition}[Twisted differential and linearized genus-$0$ complex]
% label removed: def:twisted-differential
Given $\pi\in\mathfrak g^{(0),1}$ as in Definition
\ref{def:complete-genus-filtered-package}, define
\begin{equation}
% label removed: eq:twisted-differential
d_\pi(x):=\sum_{m\ge 0}\frac1{m!}\,\ell_{m+1}^{(0)}(\pi^{\otimes m},x).
\end{equation}
This is the linearization of the genus-$0$ Maurer--Cartan equation at $\pi$.
\end{definition}

\begin{theorem}[The first modular obstruction; \ClaimStatusProvedHere]
% label removed: thm:first-modular-obstruction
Let $\mathfrak g^{\mathrm{mod}}$ be a complete genus-filtered cyclic deformation package and let
$\pi\in \mathfrak g^{(0),1}$ be its genus-$0$ Maurer--Cartan element.  Define the genus-$1$
source term
\begin{equation}
% label removed: eq:omega-one
\omega_1(\pi):=\sum_{k\ge 1}\frac1{k!}\,\ell_k^{(1)}(\pi^{\otimes k})
\in \mathfrak g^{(1),2}.
\end{equation}
Then:
\begin{enumerate}[label=\textup{(\alph*)}]
\item $d_\pi^2=0$;
\item $d_\pi\omega_1(\pi)=0$;
\item there is therefore a canonical cohomology class
\[
\Obs_1(\pi):=[\omega_1(\pi)]
\in H^2\bigl(\mathfrak g^{(1)},d_\pi\bigr);
\]
\item a genus-$1$ correction $\Theta^{(1)}\in\mathfrak g^{(1),1}$ exists modulo gauge if and
only if $\Obs_1(\pi)=0$;
\item if such genus-$1$ corrections exist, their gauge-equivalence classes form a torsor for
$H^1(\mathfrak g^{(1)},d_\pi)$.
\end{enumerate}
\end{theorem}

\begin{proof}
The linearization of the Maurer--Cartan equation at a Maurer--Cartan point always squares to
zero; this is the standard $L_\infty$ identity for the twisted unary bracket, proving $d_\pi^2=0$.
Now insert the ansatz
\[
\Pi(\h)=\pi+\h\Theta^{(1)}+O(\h^2)
\]
into \eqref{eq:full-modular-mc}.  The coefficient of $\h^0$ is the genus-$0$ Maurer--Cartan
equation for $\pi$, while the coefficient of $\h^1$ is exactly
\[
d_\pi\Theta^{(1)}+\omega_1(\pi)=0.
\]
Applying $d_\pi$ and using $d_\pi^2=0$ gives $d_\pi\omega_1(\pi)=0$, so
$\omega_1(\pi)$ defines a degree-$2$ cohomology class.  The displayed equation also shows that
a genus-$1$ lift exists precisely when this cocycle is exact.  Finally, if
$\Theta^{(1)}$ and $\widetilde\Theta^{(1)}$ are two solutions then their difference is
$d_\pi$-closed, and a genus-$1$ gauge transformation changes the difference by a
$d_\pi$-exact term.  This yields the torsor statement.
\end{proof}

\begin{definition}[Higher source terms]
% label removed: def:higher-source-terms
Let
\[
\Pi_{<G}:=\pi+\sum_{1\le h<G}\h^h\Theta^{(h)}
\]
be a partial lift satisfying the Maurer--Cartan equation modulo $\h^G$.  The
\emph{$G$-th source term} is the coefficient
\begin{equation}
% label removed: eq:higher-source-term
\omega_G(\Pi_{<G})
:=
[\h^G]\,
\sum_{r\ge 0}\sum_{k\ge 1}\frac1{k!}\,
\h^r\ell_k^{(r)}\!\bigl(\Pi_{<G}^{\otimes k}\bigr)
\in \mathfrak g^{(G),2}.
\end{equation}
\end{definition}

\begin{theorem}[All-genus obstruction tower; \ClaimStatusProvedHere]
% label removed: thm:all-genus-obstruction-tower
For every $G\ge 1$, the source term $\omega_G(\Pi_{<G})$ is $d_\pi$-closed.  Hence there is a
canonical obstruction class
\[
\Obs_G(\Pi_{<G})
:=
[\omega_G(\Pi_{<G})]
\in H^2\bigl(\mathfrak g^{(G)},d_\pi\bigr).
\]
Moreover:
\begin{enumerate}[label=\textup{(\roman*)}]
\item the partial lift $\Pi_{<G}$ extends to order $G$ if and only if
$\Obs_G(\Pi_{<G})=0$;
\item if extensions exist, their gauge-equivalence classes form a torsor for
$H^1(\mathfrak g^{(G)},d_\pi)$;
\item if $H^2(\mathfrak g^{(G)},d_\pi)=0$ for all $G\ge 1$, then $\pi$ admits a full modular
lift;
\item if, in addition, $H^1(\mathfrak g^{(G)},d_\pi)=0$ for all $G\ge 1$, then the modular lift
is unique up to gauge.
\end{enumerate}
\end{theorem}

\begin{proof}
Let
\[
E(\Pi):=\sum_{r\ge 0}\sum_{k\ge 1}\frac1{k!}\,\h^r\ell_k^{(r)}(\Pi^{\otimes k})
\]
be the Maurer--Cartan functional.  Since $\Pi_{<G}$ solves the MC equation modulo $\h^G$, one
has
\[
E(\Pi_{<G})=\h^G\omega_G(\Pi_{<G})+O(\h^{G+1}).
\]
For any $x\in\mathfrak g^{(G),1}$,
\[
E(\Pi_{<G}+\h^Gx)=\h^G\bigl(\omega_G(\Pi_{<G})+d_\pi x\bigr)+O(\h^{G+1}),
\]
because the only terms linear in $x$ of total order $G$ arise from genus-$0$ brackets with all
other entries equal to $\pi$.  Thus extension to order $G$ is equivalent to solving
\[
d_\pi x+\omega_G(\Pi_{<G})=0,
\]
which proves part \textup{(i)}.

To prove closure, apply the linearized Maurer--Cartan operator at $\Pi_{<G}$ to
$E(\Pi_{<G})$.  The $L_\infty$ identities imply the Bianchi-type identity
\[
d_{\Pi_{<G}}E(\Pi_{<G})=0.
\]
Reducing modulo $\h^{G+1}$ and using $d_{\Pi_{<G}}\equiv d_\pi \pmod{\h}$ yields
$d_\pi\omega_G(\Pi_{<G})=0$.

If $x$ and $x'$ are two order-$G$ lifts, then $x-x'$ is $d_\pi$-closed; order-$G$ gauge
transformations shift the difference by a $d_\pi$-exact term.  This gives part \textup{(ii)}.
The remaining statements follow by induction on $G$.
\end{proof}

\begin{corollary}[Existence and uniqueness by vanishing; \ClaimStatusProvedHere]
% label removed: cor:existence-uniqueness-vanishing
If
\[
H^2\bigl(\mathfrak g^{(G)},d_\pi\bigr)=0\qquad \forall G\ge 1,
\]
then every genus-$0$ Maurer--Cartan class lifts to all genera.  If moreover
\[
H^1\bigl(\mathfrak g^{(G)},d_\pi\bigr)=0\qquad \forall G\ge 1,
\]
then the lift is unique up to gauge.
\end{corollary}

\subsection{Bulk--boundary comparison as a cone problem}
% label removed: subsec:bulk-boundary-cone-problem

The obstruction tower becomes genuinely holographic only when bulk and boundary lifting
problems can be compared.  The right formal object is the cone of the linearized
bulk--boundary map.

\begin{definition}[Linearized bulk--boundary map]
% label removed: def:linearized-bulk-boundary-map
Let
\[
F:\mathfrak g^{\mathrm{mod}}_{\mathrm{bulk}}\longrightarrow \mathfrak g^{\mathrm{mod}}_{\partial}
\]
be a filtered $L_\infty$-morphism, and let
$\pi_{\mathrm{bulk}}\in\mathfrak g_{\mathrm{bulk}}^{(0),1}$,
$\pi_\partial\in\mathfrak g_\partial^{(0),1}$
be genus-$0$ Maurer--Cartan elements such that $F_*(\pi_{\mathrm{bulk}})$ is gauge-equivalent to
$\pi_\partial$.  The \emph{linearized bulk--boundary map} is the induced cochain map
\[
F_\pi:(\mathfrak g^{\mathrm{mod}}_{\mathrm{bulk}},d_{\pi_{\mathrm{bulk}}})
\longrightarrow
(\mathfrak g^{\mathrm{mod}}_{\partial},d_{\pi_\partial}).
\]
\end{definition}

\begin{theorem}[Relative bulk--boundary reconstruction; \ClaimStatusProvedHere]
% label removed: thm:relative-bulk-boundary-reconstruction
Assume the hypotheses of Definition \ref{def:linearized-bulk-boundary-map}.  If
\[
H^1\bigl(\Cone(F_\pi)\bigr)=H^2\bigl(\Cone(F_\pi)\bigr)=0,
\]
then the genus-by-genus modular lift problem for $\pi_{\mathrm{bulk}}$ is equivalent, up to gauge,
to the genus-by-genus modular lift problem for $\pi_\partial$.  In particular:
\begin{enumerate}[label=\textup{(\roman*)}]
\item every boundary modular lift uniquely determines a compatible bulk modular lift;
\item every bulk modular lift uniquely determines a compatible boundary modular lift;
\item all extension obstructions and ambiguities are measured by the same cone cohomology.
\end{enumerate}
\end{theorem}

\begin{proof}[Proof sketch]
Compatible pairs of bulk and boundary lifts form the Maurer--Cartan space of the homotopy
fiber of $F$ over $\pi_\partial$.  The tangent complex of that homotopy fiber at the point
$(\pi_{\mathrm{bulk}},\pi_\partial)$ is $\Cone(F_\pi)[-1]$.  Standard obstruction theory for formal
moduli problems therefore places the extension obstruction in
$H^2(\Cone(F_\pi))$ and the ambiguity in $H^1(\Cone(F_\pi))$.  When both groups vanish,
compatible lifts exist and are unique at each order in $\h$, hence by induction through the
entire genus filtration.
\end{proof}

\begin{corollary}[Boundary determination of the modular bulk; \ClaimStatusProvedHere]
% label removed: cor:boundary-determines-bulk
If the linearized bulk--boundary map is a quasi-isomorphism in degrees $1$ and $2$, then the
full modular bulk deformation problem is controlled by the boundary algebra.  Equivalently,
under this hypothesis the holographic dictionary is reconstructible from the boundary modular
lift alone.
\end{corollary}

\subsection{The boundary modular class}
% label removed: subsec:boundary-modular-class

The genus-$1$ source term admits a canonical graph-theoretic incarnation.  This gives the
first genuinely modular invariant of the boundary algebra.

\begin{definition}[Loop-generated modular model]
% label removed: def:loop-generated-model
A modular deformation package is \emph{loop-generated} if its first positive-genus operations
are generated entirely by connected stable graphs of first Betti number $1$ whose vertices are
decorated by genus-$0$ operations.  Equivalently, there are no primitive genus-$1$ vertices:
all genus-$1$ information is produced by a single loop with tree-level vertex decorations.
\end{definition}

\begin{definition}[One-wheel sum and boundary modular class]
% label removed: def:one-wheel-sum-and-class
Let $(A,m_\bullet)$ be a cyclic minimal $A_\infty$-algebra, and suppose its modular
deformation package is loop-generated.  Let $\mathsf{Wheel}_1$ denote the set of connected
stable graphs of first Betti number $1$, modulo isomorphism.  For
$\Gamma\in\mathsf{Wheel}_1$, let $B_\Gamma(m_\bullet)$ be the graph contraction obtained by
placing the transferred vertices $m_n$ at the vertices of $\Gamma$ and contracting internal edges
using the cyclic pairing.  Define the completed one-wheel sum
\begin{equation}
% label removed: eq:one-wheel-sum
W_1(A):=
\sum_{\Gamma\in\mathsf{Wheel}_1}^{\mathrm{comp}}
\frac{w_\Gamma}{|\Aut(\Gamma)|}
B_\Gamma(m_\bullet)
\in \mathfrak g_A^{(1),2},
\end{equation}
where $w_\Gamma$ denotes the regularized weight of the wheel type.  The cohomology class
\begin{equation}
% label removed: eq:boundary-modular-class
\mathfrak W(A):=[W_1(A)]\in H^2(\mathfrak g_A^{(1)},d_{\pi_A})
\end{equation}
is called the \emph{boundary modular class} of $A$.
\end{definition}

\begin{theorem}[One-wheel reduction; \ClaimStatusProvedHere]
% label removed: thm:one-wheel-reduction
Assume the modular deformation package of $(A,m_\bullet)$ is loop-generated.  Then:
\begin{enumerate}[label=\textup{(\roman*)}]
\item $W_1(A)$ is $d_{\pi_A}$-closed;
\item its cohomology class agrees with the first obstruction class,
\[
\mathfrak W(A)=\Obs_1(A,\pi_A);
\]
\item $\mathfrak W(A)$ depends only on the cyclic $A_\infty$ quasi-isomorphism class of $A$.
\end{enumerate}
\end{theorem}

\begin{proof}[Proof sketch]
In a loop-generated model, the coefficient of $\h$ in the Maurer--Cartan functional is a sum
over connected stable graphs of first Betti number $1$.  Each such graph has a unique loop.
Collapse every maximal tree attached to this loop by cyclic homological perturbation.  Each
collapsed tree becomes a single effective vertex labeled by one of the transferred operations
$m_n$.  The resulting expression is exactly the completed one-wheel sum \eqref{eq:one-wheel-sum}.
Therefore the order-$\h$ Maurer--Cartan equation reads
\[
d_{\pi_A}\Theta_A^{(1)}+W_1(A)=0,
\]
so $W_1(A)$ is closed and its cohomology class is the obstruction class.

Finally, cyclic $A_\infty$ quasi-isomorphisms induce quasi-isomorphisms of the corresponding
cyclic deformation complexes; the obstruction class of the Maurer--Cartan lifting problem is
therefore invariant under cyclic $A_\infty$ quasi-isomorphism.
\end{proof}

\begin{corollary}[Vanishing criterion and additivity; \ClaimStatusProvedHere]
% label removed: cor:vanishing-criterion-and-additivity
\begin{enumerate}[label=\textup{(\roman*)}]
\item A genus-$1$ modular lift exists if and only if $\mathfrak W(A)=0$.
\item If moreover $H^1(\mathfrak g_A^{(1)},d_{\pi_A})=0$, then the genus-$1$ lift is unique up to
gauge.
\item If $A=A_1\oplus A_2$ is an orthogonal direct sum of cyclic minimal $A_\infty$-algebras,
then
\[
\mathfrak W(A)=\mathfrak W(A_1)\oplus\mathfrak W(A_2).
\]
\end{enumerate}
\end{corollary}

\begin{proof}
Parts \textup{(i)} and \textup{(ii)} are immediate from Theorems
\ref{thm:first-modular-obstruction} and \ref{thm:one-wheel-reduction}.  For
\textup{(iii)}, the cyclic pairing and all transferred products split block-diagonally, so every
one-wheel contraction is supported on exactly one summand.
\end{proof}

\subsection{The holomorphic BF model on $S^3$}
% label removed: subsec:holomorphic-bf-model-on-s3

We now specialize the abstract theory to Zeng's holomorphic BF example.

\begin{definition}[The CR-cohomological boundary algebra]
% label removed: def:cr-cohomological-boundary-algebra
Let
\[
A:=\Hb = H_b^{0,0}(S^3)\oplus H_b^{0,1}(S^3).
\]
By the harmonic decomposition of the tangential Cauchy--Riemann complex and the computation
of CR cohomology, one has
\[
H_b^{0,0}(S^3)\cong \C[w_1,w_2],
\qquad
H_b^{0,1}(S^3)\cong \C[\ov w_1,\ov w_2]\,\epsilon,
\]
where $\epsilon$ is the $SU(2)$-invariant generator of $\Omega_b^{0,1}(S^3)$
\cite{ZengTwistedHolography}.  We write
\[
A^0:=\C[w_1,w_2],
\qquad
A^1:=\C[\ov w_1,\ov w_2]\,\epsilon.
\]
\end{definition}

\begin{definition}[Monomial bases and trace pairing]
% label removed: def:monomial-bases-and-trace-pairing
For $p,q,r,s\ge 0$, define
\[
a_{p,q}:=w_1^p w_2^q\in A^0,
\qquad
b_{r,s}:=\frac{(r+s+1)!}{r!\,s!}\,\ov w_1^r\ov w_2^s\,\epsilon\in A^1.
\]
The cyclic trace pairing is normalized by
\[
\Tr(a_{p,q}b_{r,s})=\delta_{pr}\delta_{qs}.
\]
\end{definition}

\begin{remark}[The cyclic SDR]
% label removed: rem:cyclic-sdr
Zeng constructs a special deformation retract
\[
h\,(\Omega_b^{0,\bullet}(S^3),\delbar_b)
\underset{i}{\overset{p}{\rightleftarrows}}
(\HCR,0)
\]
with $h\circ i=0$, $p\circ h=0$, $h^2=0$, and $\id-ip=\delbar_b h+h\delbar_b$.
Moreover the retract is compatible with the cyclic structure, and therefore transfers the
commutative dg algebra structure on $\Omega_b^{0,\bullet}(S^3)$ to a cyclic $A_\infty$ (indeed
$C_\infty$) structure $m_\bullet=\{m_n\}_{n\ge 2}$ on $A$ \cite{ZengTwistedHolography}.
This is the transferred boundary algebra governing the KK interactions.
\end{remark}

\subsection{A concrete one-vertex calculus}
% label removed: subsec:concrete-one-vertex-calculus

The first genuinely modular information enters through cyclic traces of the transferred
operations.  The one-vertex sectors exhibit a rigid arithmetic structure.

\begin{proposition}[The $m_2$ shift formula; \ClaimStatusProvedHere]
% label removed: prop:m2-shift-formula
For all $p,q,r,s\ge 0$,
\[
m_2(a_{p,q},b_{r,s})=
\begin{cases}
 b_{r-p,s-q},& p\le r,\ q\le s,\\[4pt]
 0,& \text{otherwise.}
\end{cases}
\]
\end{proposition}

\begin{proof}
Zeng's explicit formula for the transferred binary product is
\[
m_2\!\left(w_1^p w_2^q,
\frac{(r+s+1)!}{r!s!}\,\ov w_1^r \ov w_2^s\epsilon\right)
=
\frac{(r+s-p-q+1)!}{(r-p)!(s-q)!}\,\ov w_1^{r-p}\ov w_2^{s-q}\epsilon.
\]
Recognizing the right-hand side as $b_{r-p,s-q}$ proves the claim.
\end{proof}

\begin{definition}[One-vertex traced sectors]
% label removed: def:one-vertex-traced-sectors
For $n\ge 2$ and $\beta_1,\dots,\beta_{n-1}\in A^1$, define the traced one-vertex sector
\begin{equation}
% label removed: eq:one-vertex-traced-sector
Q_n(\beta_1,\dots,\beta_{n-1})
:=
\sum_{p,q\ge 0}
\Tr\!\bigl(a_{p,q}\,m_n(a_{p,q},\beta_1,\dots,\beta_{n-1})\bigr).
\end{equation}
For $n=2$, write
\[
\theta_{r,s}:=Q_2(b_{r,s}).
\]
\end{definition}

\begin{theorem}[The tadpole projector; \ClaimStatusProvedHere]
% label removed: thm:tadpole-projector
The coefficients $\theta_{r,s}$ satisfy
\[
\theta_{r,s}=
\begin{cases}
1,& r,s\ \text{both even},\\
0,& \text{otherwise}.
\end{cases}
\]
Equivalently,
\[
W^{(2)}_{\mathrm{alg}}:=\sum_{r,s\ge 0}\theta_{r,s}b_{r,s}
=\sum_{m,n\ge 0} b_{2m,2n}.
\]
In particular, the quadratic one-vertex wheel is the projector onto the even--even KK lattice.
\end{theorem}

\begin{proof}
By Proposition \ref{prop:m2-shift-formula},
\[
\theta_{r,s}
=
\sum_{p,q\ge 0}\Tr\bigl(a_{p,q} b_{r-p,s-q}\bigr)
=
\sum_{p,q\ge 0}\delta_{p,r-p}\delta_{q,s-q}.
\]
This has a solution if and only if $2p=r$ and $2q=s$, i.e. if and only if $r$ and $s$ are both
even.  When a solution exists it is unique, so the sum equals $1$.
\end{proof}

\begin{corollary}[Lowest-mode tadpole; \ClaimStatusProvedHere]
% label removed: cor:lowest-mode-tadpole
Let
\[
L:=\Span\{b_{0,0},b_{1,0},b_{0,1}\}\subset A^1.
\]
Then the tadpole projector restricts to
\[
W^{(2)}_{\mathrm{alg}}\big|_L=b_{0,0}.
\]
If one restores the geometric genus-$1$ weight $w_2$, the physical tadpole truncation on the
lowest modes is $w_2\,b_{0,0}$.
\end{corollary}

\begin{theorem}[One-vertex trace selection rule; \ClaimStatusProvedHere]
% label removed: thm:one-vertex-trace-selection-rule
Let $n\ge 2$, and let $\beta_i=b_{r_i,s_i}\in A^1$ for $1\le i\le n-1$.  Set
\[
R:=\sum_{i=1}^{n-1} r_i,
\qquad
S:=\sum_{i=1}^{n-1} s_i.
\]
Then for every $p,q\ge 0$ one has
\[
m_n(a_{p,q},b_{r_1,s_1},\dots,b_{r_{n-1},s_{n-1}})
\in
\C\cdot b_{R-p+n-2,\,S-q+n-2}.
\]
Consequently,
\[
Q_n(b_{r_1,s_1},\dots,b_{r_{n-1},s_{n-1}})\neq 0
\]
can occur only if
\begin{equation}
% label removed: eq:selection-rule
2p=R+n-2,
\qquad
2q=S+n-2.
\end{equation}
In particular, for fixed external labels there is at most one internal KK mode contributing to
the one-vertex trace.
\end{theorem}

\begin{proof}
Zeng's explicit formula for $m_n$ in the orthonormal $SU(2)$ basis shows that the output of
\[
m_n\bigl(e_{m_0}^{(j_0)},\ov e_{m_1}^{(\ov j_1)}\epsilon,\dots,
\ov e_{m_{n-1}}^{(\ov j_{n-1})}\epsilon\bigr)
\]
lies in the unique representation
\[
\ov e_{M_{n-1}}^{(\ov J_{n-1}-j_0+n-2)}\epsilon,
\qquad
\ov J_{n-1}=\sum_{i=1}^{n-1}\ov j_i,
\qquad
M_{n-1}=m_0+\sum_{i=1}^{n-1}m_i.
\]
Translate this to monomial labels by
\[
j_0=\frac{p+q}{2},\qquad m_0=\frac{p-q}{2},
\qquad
\ov j_i=\frac{r_i+s_i}{2},\qquad m_i=\frac{s_i-r_i}{2}.
\]
Writing the output as $b_{u,v}$ means
\[
\frac{u+v}{2}=\ov J_{n-1}-j_0+n-2,
\qquad
\frac{v-u}{2}=M_{n-1}.
\]
Solving for $(u,v)$ gives
\[
u=R-p+n-2,
\qquad
v=S-q+n-2.
\]
Tracing against $a_{p,q}$ forces $(u,v)=(p,q)$, which is exactly
\eqref{eq:selection-rule}.  Therefore there is at most one internal mode.
\end{proof}

\begin{theorem}[Low-mode parity law; \ClaimStatusProvedHere]
% label removed: thm:low-mode-parity-law
Restrict the one-vertex sectors to the lowest-mode space
\[
L=\Span\{b_{0,0},b_{1,0},b_{0,1}\}.
\]
For a tensor in $L^{\otimes(n-1)}$, let $N_{10}$ denote the number of $b_{1,0}$ inputs and
$N_{01}$ the number of $b_{0,1}$ inputs.  Then
\[
Q_n\big|_{L^{\otimes(n-1)}}\neq 0
\qquad\Longrightarrow\qquad
N_{10}\equiv N_{01}\equiv n\pmod 2.
\]
\end{theorem}

\begin{proof}
On $L$, one has $R=N_{10}$ and $S=N_{01}$.  By Theorem
\ref{thm:one-vertex-trace-selection-rule}, nonvanishing requires
$R+n-2$ and $S+n-2$ to be even.  Since $n-2\equiv n\pmod 2$, the statement follows.
\end{proof}

\begin{corollary}[Cubic support on the lowest modes; \ClaimStatusProvedHere]
% label removed: cor:cubic-support-lowest-modes
For $n=3$, the only nonzero low-mode one-vertex channels are the mixed ordered pairs
\[
(b_{1,0},b_{0,1})
\qquad\text{and}\qquad
(b_{0,1},b_{1,0}).
\]
No ordered pair involving $b_{0,0}$ contributes to $Q_3|_{L^{\otimes 2}}$.
\end{corollary}

\begin{proof}
For $n=3$, Theorem \ref{thm:low-mode-parity-law} requires
$N_{10}\equiv N_{01}\equiv 1\pmod 2$.  Since there are only two inputs, the only possibility is
one copy of $b_{1,0}$ and one copy of $b_{0,1}$.
\end{proof}

\begin{proposition}[Vanishing of the symmetric cubic one-vertex sector; \ClaimStatusProvedHere]
% label removed: prop:vanishing-symmetric-cubic-one-vertex-sector
The symmetric cubic one-vertex sector vanishes on the lowest modes.  Equivalently, the
trilinear form $Q_3|_{L^{\otimes 2}}$ factors through an alternating $2$-form
\[
\omega_{\mathrm{pm}}\in \Lambda^2 L^*,
\]
which we call the \emph{low-mode pre-modular two-form}.
\end{proposition}

\begin{proof}
By Corollary \ref{cor:cubic-support-lowest-modes}, the only nonzero channels are the two mixed
ordered pairs.  Zeng's antisymmetry relation for the cubic coefficients shows that swapping the
two $A^1$ inputs reverses the coefficient.  Therefore the symmetric combination cancels, and the
result factors through $\Lambda^2L^*$.
\end{proof}

\begin{theorem}[Quartic support theorem on the lowest modes; \ClaimStatusProvedHere]
% label removed: thm:quartic-support-theorem-lowest-modes
For $n=4$, the one-vertex low-mode sector is supported on exactly three unordered channels:
\[
(b_{0,0},b_{0,0},b_{0,0}),
\qquad
(b_{1,0},b_{1,0},b_{0,0}),
\qquad
(b_{0,1},b_{0,1},b_{0,0}).
\]
Equivalently, on $\Sym^3(L)$ one has a unique decomposition
\begin{equation}
% label removed: eq:q4-normal-form
Q_4\big|_{\Sym^3(L)}
=
Q_{000}\,b_{0,0}^3
+
Q_{110}\,b_{1,0}^2b_{0,0}
+
Q_{011}\,b_{0,1}^2b_{0,0},
\end{equation}
for uniquely defined scalars $Q_{000},Q_{110},Q_{011}\in\C$.
\end{theorem}

\begin{proof}
For $n=4$, Theorem \ref{thm:low-mode-parity-law} gives
$N_{10}\equiv N_{01}\equiv 0\pmod 2$.  Since there are three inputs,
\[
N_{10}+N_{01}+N_{00}=3,
\]
so the only possible triples are
\[
(N_{10},N_{01},N_{00})=(0,0,3),\ (2,0,1),\ (0,2,1).
\]
These are exactly the three unordered channels above.
\end{proof}

\begin{corollary}[The charged quartic sector is two-dimensional; \ClaimStatusProvedHere]
% label removed: cor:charged-quartic-sector-two-dimensional
After removing the neutral channel $b_{0,0}^3$, the first charged one-vertex quartic sector on
low modes lies in the $2$-dimensional space
\[
\Span\{b_{1,0}^2b_{0,0},\ b_{0,1}^2b_{0,0}\}.
\]
Thus the first charged quartic modular datum is encoded by the pair $(Q_{110},Q_{011})$.
\end{corollary}

\begin{remark}[Tiny internal support]
% label removed: rem:tiny-internal-support
The trace selection rule determines the internal test mode uniquely.  For the three quartic
channels of Theorem \ref{thm:quartic-support-theorem-lowest-modes}, one finds
\[
Q_{000}: \ a_{1,1},
\qquad
Q_{110}: \ a_{2,1},
\qquad
Q_{011}: \ a_{1,2}.
\]
So the charged quartic coefficients are algebraically finite before any regularization issue can
arise.  This contrasts sharply with the generic loop sectors of the KK theory,
where infinite KK sums do appear \cite{ZengTwistedHolography}.
\end{remark}

\subsection{The first two-vertex bubble and the exceptional finite locus}
% label removed: subsec:first-two-vertex-bubble-exceptional-finite-locus

The first genuinely loop-shaped graph is the two-vertex cubic--cubic wheel.  Even here, the
lowest-mode sector remains exceptionally rigid.

\begin{theorem}[Support of the first mixed cubic--cubic bubble; \ClaimStatusProvedHere]
% label removed: thm:support-first-mixed-cubic-cubic-bubble
Consider the first two-vertex wheel built from two cubic vertices.  Restricted to the lowest-mode
subspace $L$, its symmetric external content is supported on a single monomial:
\[
b_{1,0}^2 b_{0,1}^2.
\]
No monomial involving $b_{0,0}$ occurs in this first cubic--cubic bubble sector.
\end{theorem}

\begin{proof}
By Corollary \ref{cor:cubic-support-lowest-modes}, each cubic vertex is nonzero only if its two
$A^1$ inputs are one copy of $b_{1,0}$ and one copy of $b_{0,1}$.  Therefore every nonzero
cubic--cubic wheel contributes exactly one $b_{1,0}$ and one $b_{0,1}$ from each vertex.  The
total external monomial is thus forced to be $b_{1,0}^2b_{0,1}^2$.
\end{proof}

\begin{theorem}[Exceptional finiteness of the first mixed bubble; \ClaimStatusProvedHere]
% label removed: thm:exceptional-finiteness-first-mixed-bubble
The first mixed cubic--cubic bubble on the lowest modes is finite before regularization.  More
precisely, for the mixed cubic sector one has
\[
m_3(a_{p,q},b_{1,0},b_{0,1})\in \C\cdot b_{2-p,2-q},
\]
and therefore only the finite range
\[
0\le p,q\le 2
\]
can contribute to the first mixed bubble coefficient.  In particular, the first mixed bubble is a
finite algebraic contraction of finitely many cubic structure constants.
\end{theorem}

\begin{proof}
Apply Theorem \ref{thm:one-vertex-trace-selection-rule} with $n=3$ and external labels
$(r_1,s_1)=(1,0)$, $(r_2,s_2)=(0,1)$.  Then $R=S=1$, so the cubic output is
\[
b_{R-p+1,S-q+1}=b_{2-p,2-q}.
\]
For this to make sense in $A^1$, one must have $2-p\ge 0$ and $2-q\ge 0$, hence
$0\le p,q\le 2$.  Thus the internal mode sum is finite.
\end{proof}

\begin{definition}[The low-mode charged modular jet]
% label removed: def:low-mode-charged-modular-jet
The \emph{low-mode charged modular jet} of the holomorphic BF boundary algebra is the triple
\begin{equation}
% label removed: eq:charged-modular-jet
\mathfrak J_{\mathrm{ch}}^{\mathrm{low}}:=(Q_{110},Q_{011},B_{\mathrm{mix}}),
\end{equation}
where $Q_{110}$ and $Q_{011}$ are the charged one-vertex quartic coefficients of
\eqref{eq:q4-normal-form}, and $B_{\mathrm{mix}}$ is the coefficient of
$b_{1,0}^2b_{0,1}^2$ in the first cubic--cubic bubble sector.
\end{definition}

\begin{theorem}[Low-mode modular normal form; \ClaimStatusProvedHere]
% label removed: thm:low-mode-modular-normal-form
Up to and including the first cubic--cubic bubble order, the symmetric modular shadow of the
holomorphic BF boundary algebra on the lowest-mode sector has the form
\begin{equation}
% label removed: eq:low-mode-modular-normal-form
W_{\mathrm{low}}
=
w_2\,b_{0,0}
+
Q_{000}\,b_{0,0}^3
+
Q_{110}\,b_{1,0}^2b_{0,0}
+
Q_{011}\,b_{0,1}^2b_{0,0}
+
B_{\mathrm{mix}}\,b_{1,0}^2b_{0,1}^2
+
\textup{(higher sectors)}.
\end{equation}
No other symmetric low-mode monomial can appear at this order.
\end{theorem}

\begin{proof}
The first term is the tadpole contribution from Corollary
\ref{cor:lowest-mode-tadpole}.  The one-vertex cubic symmetric contribution vanishes by
Proposition \ref{prop:vanishing-symmetric-cubic-one-vertex-sector}.  The one-vertex quartic
contribution is exhausted by Theorem \ref{thm:quartic-support-theorem-lowest-modes}.  The
first two-vertex bubble contribution is exhausted by Theorem
\ref{thm:support-first-mixed-cubic-cubic-bubble}.  No further symmetric low-mode monomial can
occur before higher-valence or higher-vertex sectors enter.
\end{proof}

\begin{remark}[What has been proved]
% label removed: rem:what-has-been-proved
Theorem \ref{thm:low-mode-modular-normal-form} is the sharpest rigorous reduction obtained
in this section.  It does \emph{not} compute the three charged numbers
$Q_{110},Q_{011},B_{\mathrm{mix}}$, but it proves that these are the \emph{only} charged low-mode
symmetric modular coefficients that can appear at the first nontrivial orders.  This converts an
infinite KK/modular problem into a finite jet problem.
\end{remark}

\subsection{A strengthened frontier: new structures and new targets}
% label removed: subsec:strengthened-frontier-new-structures

The reduction above produces several new mathematical structures.

\begin{definition}[The low-mode pre-modular package]
% label removed: def:low-mode-pre-modular-package
The \emph{low-mode pre-modular package} of the holomorphic BF boundary algebra is the pair
\[
\bigl(\omega_{\mathrm{pm}},\mathfrak J_{\mathrm{ch}}^{\mathrm{low}}\bigr),
\]
where $\omega_{\mathrm{pm}}\in \Lambda^2L^*$ is the alternating cubic one-vertex shadow of
Proposition \ref{prop:vanishing-symmetric-cubic-one-vertex-sector}, and
$\mathfrak J_{\mathrm{ch}}^{\mathrm{low}}$ is the charged modular jet of Definition
\ref{def:low-mode-charged-modular-jet}.
\end{definition}

\begin{proposition}[Functoriality under cyclic equivalence on the lowest modes; \ClaimStatusProvedHere]
% label removed: prop:functoriality-lowest-modes
Let $A$ and $A'$ be cyclic minimal $A_\infty$ models whose lowest-mode sectors are identified
by a cyclic $A_\infty$ quasi-isomorphism preserving the subspace
$L=\Span\{b_{0,0},b_{1,0},b_{0,1}\}$.  Then:
\begin{enumerate}[label=\textup{(\roman*)}]
\item the low-mode pre-modular two-form $\omega_{\mathrm{pm}}$ is carried to the corresponding
$\omega_{\mathrm{pm}}'$;
\item the support constraints of Theorems
\ref{thm:quartic-support-theorem-lowest-modes} and
\ref{thm:support-first-mixed-cubic-cubic-bubble} are preserved;
\item the charged modular jet transforms functorially.
\end{enumerate}
\end{proposition}

\begin{proof}
The support theorems depend only on the transferred $A_\infty$ structure, the cyclic pairing,
and the grading/weight bookkeeping encoded by the selection rule.  A cyclic
$A_\infty$ quasi-isomorphism preserving $L$ transports all of this data.  Therefore the induced
trilinear and quartic forms on $L$ and the first bubble coefficient transform functorially.
\end{proof}

\begin{problem}[The quartic Wigner identity problem]
% label removed: prob:quartic-wigner-identity
Find a closed-form representation-theoretic identity, directly in terms of Clebsch--Gordan and
Wigner-symbol data, for the charged quartic coefficients
\[
Q_{110}=\Tr\bigl(a_{2,1}\,m_4(a_{2,1},b_{1,0},b_{1,0},b_{0,0})\bigr),
\qquad
Q_{011}=\Tr\bigl(a_{1,2}\,m_4(a_{1,2},b_{0,1},b_{0,1},b_{0,0})\bigr).
\]
The support theory developed above shows that these are finite, tiny-support Wigner sums.
The missing theorem is a conceptual closed-form evaluation.
\end{problem}

\begin{problem}[Universal finite-locus theorem]
% label removed: prob:universal-finite-locus
Characterize, in an arbitrary loop-generated cyclic $A_\infty$ model, exactly which low-mode
wheel sectors remain finite before regularization.  The holomorphic BF calculation suggests the
existence of an \emph{exceptional finite locus}: a representation-theoretically defined region of
external weight space on which loop diagrams collapse to finite algebraic contractions even
though generic loop sectors diverge.
\end{problem}

\begin{problem}[The modular bootstrap problem]
% label removed: prob:modular-bootstrap
Use the quadratic identities of the BRST anomaly/higher-operation formalism to bootstrap the
first genuinely modular boundary coefficients from associativity and stable-graph factorization,
rather than direct Feynman summation.  The expected endpoint is a boundary modular analog of
the bootstrap constraints on celestial one-loop OPE data.
\end{problem}

\begin{problem}[Quantization of the boundary modular class]
% label removed: prob:quantization-boundary-modular-class
Relate the genus-$1$ boundary modular class and the charged modular jet to deformation
quantization of the underlying coisson/Poisson-vertex structure.  The natural bridge is the
$3$d holomorphic--topological Poisson-vertex framework of Khan--Zeng, together with the
configuration-space/Feynman realization of higher operations and their quadratic identities.
\end{problem}

\subsection{Thread-complete provisional data package}
% label removed: subsec:thread-complete-provisional-data-package

This subsection records, in a quarantined form, the exact low-mode coefficient identities that
were discussed in the working thread.  They are included \emph{only} to preserve completeness of
the development.  They are \emph{not} used anywhere else in this file, and they should be treated
as provisional until independently verified by a direct symbolic derivation from Zeng's explicit
$A_\infty$ formulas.

\begin{conjecture}[Low mixed cubic reflection law, provisional; \ClaimStatusConjectured]
% label removed: conj:low-mixed-cubic-reflection-law-provisional
Let
\[
A:=b_{1,0},
\qquad
B:=b_{0,1}.
\]
Suppose the low mixed cubic coefficients are defined by
\[
m_3(a_{p,q},A,B)=c^{AB}_{p,q}\,b_{2-p,2-q},
\qquad
m_3(a_{p,q},B,A)=c^{BA}_{p,q}\,b_{2-p,2-q},
\qquad 0\le p,q\le 2.
\]
Then one expects the exact reflection identity
\[
c^{BA}_{p,q}=-c^{AB}_{2-p,2-q}.
\]
This should be the sharp low-mode shadow of Zeng's antisymmetry relation for $m_3$.
\end{conjecture}

\begin{conjecture}[Explicit charged quartic jet, provisional; \ClaimStatusConjectured]
% label removed: conj:explicit-charged-quartic-jet-provisional
The ordered charged quartic coefficients satisfy
\[
Q_{110}=Q_{011}=-\frac1{15},
\]
and after symmetrization in the boundary OPE one obtains
\[
Q_{110}^{\mathrm{sym}}=Q_{011}^{\mathrm{sym}}=\frac1{18}.
\]
\end{conjecture}

\begin{conjecture}[First mixed bubble coefficient, provisional; \ClaimStatusConjectured]
% label removed: conj:first-mixed-bubble-coefficient-provisional
The algebraic first mixed cubic--cubic bubble coefficient is
\[
B_{\mathrm{mix}}^{\mathrm{alg}}=-\frac{11}{9}.
\]
Equivalently, if $w_{3,3}$ denotes the geometric/Feynman weight of the $(3,3)$ wheel type,
then the full coefficient is expected to be
\[
B_{\mathrm{mix}}=-\frac{11}{9}\,w_{3,3}.
\]
\end{conjecture}

\begin{remark}[How the provisional package should be used]
% label removed: rem:how-to-use-provisional-package
The rigorous content of the present section is the reduction to the three-parameter jet
$(Q_{110},Q_{011},B_{\mathrm{mix}})$ together with the vanishing/support/finiteness theorems that
force this reduction.  The three exact rational values recorded in Conjectures
\ref{conj:explicit-charged-quartic-jet-provisional} and
\ref{conj:first-mixed-bubble-coefficient-provisional} should therefore be viewed as
\emph{targets} for the next theorem-sized computation, not as premises for further arguments.
\end{remark}

\begin{proposition}[Cross-conjecture relation; \ClaimStatusProvedHere]
% label removed: prop:celestial-cross-conjecture
The three provisional conjectures are not independent.
The ordered-to-symmetrized correction satisfies the structural identity
\begin{equation}% label removed: eq:cross-conjecture-structural
Q^{\mathrm{sym}}_{110} - Q^{\mathrm{ord}}_{110} \;=\; -\frac{B_{\mathrm{mix}}}{10}.
\end{equation}
If
Conjectures~\textup{\ref{conj:explicit-charged-quartic-jet-provisional}}
and~\textup{\ref{conj:first-mixed-bubble-coefficient-provisional}} hold,
then
$Q^{\mathrm{sym}}_{110} - Q^{\mathrm{ord}}_{110} = 11/90$.
\end{proposition}

\begin{proof}
The symmetrization $Q^{\mathrm{sym}} - Q^{\mathrm{ord}}$ at
arity~$4$ receives a single correction from the cubic--cubic
bubble graph: the $(3,3)$ wheel with mixed insertions.
By direct expansion of the symmetrization formula for~$m_4$
in terms of~$m_3 \circ m_3$ contractions, each such bubble
contributes with coefficient~$-1/10$ (the inverse of the
number of ordered placements of two cubic vertices on a
$4$-input tree).
Hence
$Q^{\mathrm{sym}}_{110} - Q^{\mathrm{ord}}_{110}
= -B_{\mathrm{mix}}/10$.
The numerical value $11/90$ follows from
$B_{\mathrm{mix}}^{\mathrm{alg}} = -11/9$
(Conjecture~\ref{conj:first-mixed-bubble-coefficient-provisional}).
\end{proof}

\subsection{The distilled picture}
% label removed: subsec:distilled-picture

The entire thread can now be summarized in a single hierarchy.

\begin{enumerate}[label=\textup{(\arabic*)}]
\item \textbf{Disk level.}
The KK reduction of Zeng produces a cyclic transferred $A_\infty$ algebra on
$\Hb$, and the boundary OPE of the holomorphic BF model is built from the operations
$m_2,m_3,m_4,\dots$.

\item \textbf{Modular level.}
The correct higher object is a modular Maurer--Cartan lift of the disk-level class.  The first
obstruction is the genus-$1$ class $\Obs_1$, and in loop-generated models it is represented by
the one-wheel class $\mathfrak W(A)$.

\item \textbf{Holographic level.}
Twisted/celestial holography becomes the assertion that bulk and boundary modular lifting
problems coincide.

\item \textbf{Concrete BF level.}
On the lowest KK sector, the first symmetric modular problem collapses to the finite normal
form \eqref{eq:low-mode-modular-normal-form}.  Hence the first charged symmetric modular
information is controlled by the single jet
\[
\mathfrak J_{\mathrm{ch}}^{\mathrm{low}}=(Q_{110},Q_{011},B_{\mathrm{mix}}).
\]
\end{enumerate}

This is the strengthened first-principles form of the result.  The disk-level boundary
VOA is not the end of the theory.  It is the shadow of a modular Maurer--Cartan problem; the
first modular class is a new invariant of the boundary algebra; and in the holomorphic BF test
case the infinite KK problem reduces, at low modes and low modular order, to a finite and
rigid charged jet.

%% ===================================================================
%%  THE MODULAR HOLOGRAPHY PROGRAMME
%% ===================================================================

\providecommand{\ModEnv}{\operatorname{ModEnv}}
\providecommand{\Fact}{\operatorname{Fact}}
\providecommand{\ot}{\widehat\otimes}

\section{The modular holography programme}
% label removed: sec:modular-holography-programme

The preceding sections built the abstract modular obstruction tower and tested it on a single
example.  We now step back and formulate the broader programme that these constructions
point toward: a systematic passage from genus-zero holographic data to a full modular
homotopy type.

\subsection{The genus-zero holographic package (synthesis)}
% label removed: subsec:genus-zero-holographic-package

The first step is to collect the genus-zero data established by the recent perturbative
literature into a single coherent package.

\begin{theorem}[Genus-zero holographic package; \ClaimStatusProvedElsewhere]
% label removed: thm:genus-zero-holographic-package
For every perturbative holomorphic--topological theory $T$ of the type considered in
\cite{CDG20, Zeng23, GKW24, KZ25, DNP25}, the genus-zero datum
\[
\mathcal G_0(T;B)
\;=\;
\bigl(V_T,\;\{-,-\}_T,\;V_{\partial,T}(B),\;\beta_B,\;A_T^!,\;C_T,\;r_T(z),\;\{m_k^T\}_{k\ge 2}\bigr)
\]
is defined, where:
\begin{enumerate}[label=\textup{(\arabic*)}]
\item $V_T$ is a commutative chiral algebra equipped with a cohomological degree $-1$ Poisson
bracket $\{-,-\}_T$.  For $3$d $\mathcal N=2$ holomorphic twists, there is a secondary stress
tensor.

\item $V_{\partial,T}(B)$ is a boundary chiral algebra, depending on a transverse boundary
condition $B$, together with a bulk--boundary map
\[
\beta_B:V_T\longrightarrow Z\bigl(V_{\partial,T}(B)\bigr)
\]
into the centre of the boundary algebra.

\item In Kaluza--Klein reduction one can choose $B=B_{\mathrm{univ}}$ so that the boundary algebra
coincides with the universal defect chiral algebra of the unreduced model.

\item Deformation, anomaly, and generalized OPE data are controlled by higher operations
$\{m_k^T\}_{k\ge 2}$ satisfying quadratic (Wess--Zumino) identities.

\item Perturbative line operators are $A_T^!$-modules; their OPE is controlled by a
Maurer--Cartan kernel
\[
r_T(z)\;\in\;A_T^!\;\ot\;A_T^!\bigl(\!(z^{-1})\!\bigr)
\]
satisfying an $A_\infty$ Yang--Baxter constraint (dg-shifted Yangian).
\end{enumerate}
\end{theorem}

\begin{proof}[Sources]
Item~\textup{(1)} and item~\textup{(2)} are established in
\cite{CDG20}.  Item~\textup{(3)} is the main theorem of \cite{Zeng23}.  Item~\textup{(4)}
follows from the higher-operation formalism of \cite{GKW24}.  Item~\textup{(5)} is the
principal construction of \cite{DNP25}.
\end{proof}

\begin{remark}
% label removed: rem:genus-zero-package-modular-context
The genus-zero package theorem is not a single new result; it is a deliberate assembly of five
strands of the recent literature into one formal object $\mathcal G_0(T;B)$.  The components listed above are exactly the data that must be completed to higher genus.
The next two definitions make that completion precise.
\end{remark}

\subsection{The full modular homotopy type}
% label removed: subsec:full-modular-homotopy-type

\begin{definition}[Full modular homotopy type; \ClaimStatusConjectured]
% label removed: def:full-modular-homotopy-type
A \emph{full modular homotopy type} for $T$ based at $B$ is a completion of the genus-zero
package $\mathcal G_0(T;B)$ by multilinear maps
\begin{equation}
% label removed: eq:modular-operations
\mu_{g,n}^T:\;
C_\bullet(\ov{\M}_{g,n+1})\;\otimes\;(A_T^!)^{\ot\, n}
\;\longrightarrow\;A_T^!,
\qquad 2g-2+n>0,
\end{equation}
subject to the following axioms.
\begin{enumerate}[label=\textup{(\alph*)}]
\item \emph{Genus-zero recovery.}  The operations $\mu_{0,n}^T$ recover the local higher
operations $m_n^T$ of the genus-zero package.

\item \emph{Stable-graph gluing.}  The family $\{\mu_{g,n}^T\}$ is compatible with the
stable-graph composition law on the modular operad
$C_\bullet(\ov{\M}_{\bullet,\bullet})$ in the sense of Getzler--Kapranov \cite{GeK98}.

\item \emph{Line-kernel refinement.}  The genus-zero Maurer--Cartan kernel $r_T(z)$ extends to
a formal genus expansion
\begin{equation}
% label removed: eq:modular-r-matrix-celestial
R_T^{\mathrm{mod}}(z;\h)
\;=\;\sum_{g\ge 0}\h^{2g}\,r_{T,g}(z),
\end{equation}
with $r_{T,0}=r_T$ and
\[
r_{T,g}(z)\;\in\;
H^\bullet(\ov{\M}_{g,2})\;\otimes\;A_T^!\;\ot\;A_T^!\bigl(\!(z^{-1})\!\bigr).
\]
\end{enumerate}
\end{definition}

\begin{definition}[Modular state spaces and partition function]
% label removed: def:modular-state-spaces-partition-function
Given a full modular homotopy type for $T$ based at $B$ and line operators
$L_1,\dots,L_n$, define the \emph{modular state space}
\begin{equation}
% label removed: eq:modular-state-space
H_{g,n}^{T;B}(L_1,\dots,L_n)
\;:=\;
\int_{\Sigma_{g,n}}\Fact_{T;B}[L_1,\dots,L_n],
\end{equation}
where the right-hand side denotes factorization (chiral) homology of the factorization algebra
$\Fact_{T;B}$ over the stable pointed curve $\Sigma_{g,n}$, in the sense of Ayala--Francis
\cite{AF15} and Beilinson--Drinfeld \cite{BD04}.  The \emph{partition function} is the
generating series
\begin{equation}
% label removed: eq:partition-function
Z_{T;B}(\h,t)
\;:=\;
\sum_{\substack{g,n\ge 0\\2g-2+n>0}}
\h^{2g-2+n}\,\frac{1}{n!}\,\chi_{\gr}\!\bigl(H_{g,n}^{T;B}(t)\bigr),
\end{equation}
where $\chi_{\gr}$ is the graded Euler characteristic.
\end{definition}

\subsection{The modular envelope conjecture}
% label removed: subsec:modular-envelope-conjecture

\begin{conjecture}[Modular envelope; \ClaimStatusConjectured]
% label removed: conj:modular-envelope
For the classes of theories appearing in Theorem~\ref{thm:genus-zero-holographic-package},
every genus-zero package $\mathcal G_0(T;B)$ admits a full modular homotopy type
$\ModEnv(T;B)$.
\end{conjecture}

\begin{remark}
% label removed: rem:modular-envelope-strength
Conjecture~\ref{conj:modular-envelope} is deliberately stronger than the mere existence of
genus-by-genus correlation functions.  It asks for an actual modular operadic object whose
gluing law controls higher operations, line OPE, and defect observables in a single package.
The factorization-homology viewpoint of Definition
\ref{def:modular-state-spaces-partition-function} is natural for this purpose because it
integrates local algebraic data over a manifold and comes equipped with local-to-global
assembly, excision, and Poincar\'e/Koszul duality features \cite{AF15}.  The modular envelope
conjecture therefore sits at the intersection of the abstract modular lifting problem
(Conjecture~\ref{conj:modular-twisted-celestial}) and the operadic completion programme of
Getzler--Kapranov \cite{GeK98}.

The obstructions to existence of $\ModEnv(T;B)$ are deformation-theoretic: at each genus~$g$,
the extension from genus $\leq g-1$ to genus~$g$ is obstructed by a class in~$H^2$ of the
modular cyclic deformation complex $\mathrm{Def}_{\mathrm{cyc}}^{\mathrm{mod}}(A_T)$
(Definition~\ref*{def:modular-cyclic-deformation-complex} of Volume~I).  In the free case
($\beta\gamma$ system), all higher obstructions vanish and the envelope is determined by
determinant-line combinatorics.  For non-free theories ($W_3$, KS holography), the obstructions
are controlled by the shadow obstruction tower: the arity-$r$ shadow
$\mathrm{Sh}_r(A_T)$ governs the obstruction at $r$-point level.
\end{remark}

\subsection{The research programme}
% label removed: subsec:modular-holography-research-programme

The programme can be summarized: build modular homotopy types by
combining stable-graph summation with logarithmic configuration geometry
and dg-shifted Yangian line algebras.  Three guiding questions organize
the programme.

\begin{problem}[Existence]
% label removed: prob:modular-envelope-existence
Can $\ModEnv(T;B)$ be constructed from the genus-zero package $\mathcal G_0(T;B)$ by
bar-cobar or deformation-theoretic completion over the modular operad
$C_\bullet(\ov{\M}_{g,n})$?  In the language of this monograph, the question is whether the
modular lift of Section~\ref{subsec:order-by-order-obstruction-theory} can be upgraded from
an order-by-order obstruction-theory statement to an explicit algebraic construction via
the bar-cobar machinery of Part~I.
\end{problem}

\begin{problem}[Geometry]
% label removed: prob:modular-envelope-geometry
Can the gluing operations $\mu_{g,n}^T$ be realized by explicit integrals over logarithmic
Fulton--MacPherson spaces \cite{Mok25} and their degeneration correspondences?  More
precisely, does the stable-graph composition law of
Definition~\ref{def:full-modular-homotopy-type}\textup{(b)} admit a geometric incarnation in
which vertex insertions are integrated over logarithmic configuration strata and edge
propagators are integrated over degeneration divisors?
\end{problem}

\begin{problem}[Computation]
% label removed: prob:modular-envelope-computation
Can one compute the first genuinely modular coefficients of
$R_T^{\mathrm{mod}}(z;\h)$ in the basic examples?  The holomorphic BF model
(Section~\ref{subsec:holomorphic-bf-model-on-s3}) provides the natural test case: the finite
charged jet $\mathfrak J_{\mathrm{ch}}^{\mathrm{low}}$ of
Definition~\ref{def:low-mode-charged-modular-jet} should be recoverable from the genus-one
coefficient $r_{T,1}(z)$ of the formal ordered arity-$2$ shadow series
\eqref{eq:modular-r-matrix}, restricted to lowest KK modes.
\end{problem}

%%%%%%%%%%%%%%%%%%%%%%%%%%%%%%%%%%%%%%%%%%%%%%%%%%%%%%%%%%%%%%%%%%%%%%%%%%%%%
% WORKED EXAMPLES: THE THREE HOLOGRAPHIC ATOMS
%%%%%%%%%%%%%%%%%%%%%%%%%%%%%%%%%%%%%%%%%%%%%%%%%%%%%%%%%%%%%%%%%%%%%%%%%%%%%
\section{Worked examples: the three holographic atoms}
% label removed: sec:three-holographic-atoms

We now instantiate the modular holography programme on three representative examples of
increasing complexity.  Each illustrates a distinct stratum of the machine: the free
$\beta\gamma$ system tests the formalism at the strictest possible level, Kodaira--Spencer
holography tests the open/closed interface with an explicit universal defect algebra, and
the $\mathcal{W}_3$ model provides the first nonlinear higher-spin arena.

\subsection{The free \texorpdfstring{$\beta\gamma$}{betagamma} holographic atom}
% label removed: subsec:betagamma-holographic-atom

The simplest holographic atom is the free $\beta\gamma$ system.

\begin{definition}[Free $\beta\gamma$ chiral algebra]
% label removed: def:betagamma-chiral-algebra
Let $V_{\beta\gamma} := \Sym\bigl(\C[\partial]\beta \oplus \C[\partial]\gamma\bigr)$
with $\lambda$-bracket $\{\beta_\lambda\,\gamma\} = 1$.  The OPE is
$\beta(z)\gamma(w) \sim (z-w)^{-1}$; all higher products vanish.
\end{definition}

The 3d origin is Costello--Dimofte--Gaiotto's interval compactification \cite{CDG20}:
a 3d $\N=2$ free chiral multiplet on $\C \times [0,1]$ with fields $\phi(z,t)$ and
$\psi^{(1)}(z,t)$.  Setting
\begin{equation}
% label removed: eq:betagamma-interval
\beta(z) := \phi(z,0),
\qquad
\gamma(z) := \int_0^1 \psi^{(1)}(z,t)\,dt,
\end{equation}
and applying the boundary conditions produces the full $\beta\gamma$ system as the
boundary chiral algebra.

\begin{proposition}[Minimal modular content; \ClaimStatusProvedElsewhere]
% label removed: prop:betagamma-modular-content
The free $\beta\gamma$ system is a minimal modular candidate in the following sense:
\begin{enumerate}[label=\textup{(\roman*)}]
\item The genus-zero Maurer--Cartan kernel $r_{\beta\gamma}(z)$ is the standard propagator
  $(z-w)^{-1}$;
\item the modular operation $\mu_{0,2} \neq 0$ (the OPE itself), but all higher local
  operations $m_{k \geq 3}$ vanish identically;
\item genus dependence enters exclusively through factorization homology over stable curves.
\end{enumerate}
\end{proposition}

The modular state spaces take the form
\begin{equation}
% label removed: eq:betagamma-state-spaces
H_{g,n}^{\beta\gamma}(L_1,\dots,L_n)
\;=\;
\int_{\Sigma_{g,n}} \Fact_{\beta\gamma}[L_1,\dots,L_n],
\end{equation}
where the factorization algebra $\Fact_{\beta\gamma}$ is freely generated.  This is the
cleanest test case: the free OPE ensures that all modular complexity resides in the
determinant-line genus dependence and the topology of $\ov{\M}_{g,n}$.

\subsection{Kodaira--Spencer holography and the open-closed modular package}
% label removed: subsec:kodaira-spencer-holography

Kodaira--Spencer holography provides a richer test, with an explicit universal defect
algebra realizing the open-closed decomposition.

\begin{definition}[Zeng's boundary algebra; \ClaimStatusProvedElsewhere]
% label removed: def:zeng-boundary-algebra
Following Zeng \cite{Zeng23}, the boundary chiral algebra for Kodaira--Spencer gravity
in the large-$N$ limit is constructed from:
\begin{itemize}
\item a $(b,c)$ ghost system in $\mathfrak{gl}_N$,
\item adjoint $\beta\gamma$ pairs $(Z_1,Z_2)$, and
\item fundamental $\beta\gamma$ pairs $(I,J)$,
\end{itemize}
with BRST charge
\begin{equation}
% label removed: eq:ks-brst
Q_{\mathrm{BRST}}
\;=\;
\oint\!\Bigl(\Tr(c\,Z_1\,Z_2) + \Tr(I\,c\,J) + \tfrac{1}{2}\Tr(b\,c^2)\Bigr).
\end{equation}
\end{definition}

\begin{proposition}[Single-trace BRST cohomology; \ClaimStatusProvedElsewhere]
% label removed: prop:ks-single-trace
The large-$N$ BRST cohomology of the Kodaira--Spencer boundary algebra decomposes into
single-trace towers:
\[
A(n),\quad B(n),\quad C(n),\quad D(n),\quad E_t(n),
\qquad n \geq 1.
\]
The first four towers $A$, $B$, $C$, $D$ match the single-trace operators of
Kodaira--Spencer gravity.  The $E$-tower matches holomorphic Chern--Simons theory.
\end{proposition}

\begin{definition}[Universal defect algebra]
% label removed: def:universal-defect-algebra-ks
The \emph{universal defect algebra} is
\begin{equation}
% label removed: eq:universal-defect-ks
U_{\mathrm{KS}} := A_{\partial,\mathrm{KS}}(B_{\mathrm{univ}}),
\end{equation}
with \emph{closed-string sector}
\[
U_{\mathrm{KS}}^{\mathrm{cl}}
\;:=\;
\bigl\langle A(n),\;B(n),\;C(n),\;D(n) \bigr\rangle.
\]
The modular state spaces and partition function are
\begin{align}
% label removed: eq:ks-state-spaces
H_{g,n}^{\mathrm{KS}}
&\;:=\;
\int_{\Sigma_{g,n}} U_{\mathrm{KS}},
&
H_{g,n}^{\mathrm{cl,KS}}
&\;:=\;
\int_{\Sigma_{g,n}} U_{\mathrm{KS}}^{\mathrm{cl}},
\\
% label removed: eq:ks-partition-function
Z_{\mathrm{KS}}(\h)
&\;:=\;
\sum_{g \geq 0} \h^{2g-2}\,\chi_{\gr}\!\bigl(H_g^{\mathrm{KS}}\bigr).
\end{align}
\end{definition}

\begin{conjecture}[Open-closed modular splitting; \ClaimStatusConjectured]
% label removed: conj:ks-open-closed-splitting-celestial
There is a splitting of modular state spaces
\begin{equation}
% label removed: eq:ks-open-closed-splitting
H_{g,n}^{\mathrm{KS}}
\;\simeq\;
H_{g,n}^{\mathrm{cl,KS}}
\;\hat\otimes\;
H_{g,n}^{\mathrm{gauge,KS}},
\end{equation}
where the gauge factor $H_{g,n}^{\mathrm{gauge,KS}}$ is controlled by the $E$-tower
and carries a filtration compatible with the 't~Hooft expansion.
\end{conjecture}

\subsection{The \texorpdfstring{$\mathcal{W}_3$}{W3} modular envelope}
% label removed: subsec:w3-modular-envelope

The $\mathcal{W}_3$ algebra is the simplest nonlinear example and the first
higher-spin test of the full modular machine.

\begin{definition}[Classical $\mathcal{W}_3$ PVA]
% label removed: def:classical-w3-pva
The classical $\mathcal{W}_3$ Poisson vertex algebra is
$P_{\mathcal{W}_3} := \Sym\bigl(\C[\partial]T \oplus \C[\partial]W\bigr)$,
with $\lambda$-brackets determined by conformal weights $(2,3)$ and the standard
Zamolodchikov normalization.
\end{definition}

The 3d origin is the holomorphic-topological Poisson sigma model of Khan--Zeng
\cite{KZ25}, with target space determined by $P_{\mathcal{W}_3}$ and fields
$(T,\mu)$ (stress tensor and Beltrami differential) together with $(W,\chi)$
(spin-3 current and spin-3 source).  The topological enhancement comes from the
Virasoro element $T$, which provides a secondary stress tensor.

\begin{conjecture}[Slab reduction; \ClaimStatusConjectured]
% label removed: prop:w3-slab-reduction
The slab reduction of the $\mathcal{W}_3$ Poisson sigma model produces effective
central charges:
\[
c^{(2)} = 26,
\qquad
c^{(3)} = 74,
\qquad
c_{\mathrm{pair}} = 100,
\qquad
c_{\mathrm{eff}} = c + 50.
\]
\end{conjecture}

\begin{definition}[Modular $\mathcal{W}_3$ envelope; \ClaimStatusConjectured]
% label removed: def:modular-w3-envelope
Assuming the conjectural slab-reduction package above, the
\emph{expected modular $\mathcal{W}_3$ envelope} is the tuple
\begin{equation}
% label removed: eq:modular-w3-envelope
\ModEnv(\mathcal{W}_3)
\;=\;
\bigl(\mathcal{W}_3^q(c_{\mathrm{eff}}),\;
(\mathcal{W}_3^q)^!,\;
Y_{\mathcal{W}_3}^{\mathrm{mod}},\;
\Theta_{g,n}^{\mathcal{W}_3},\;
F_{\mathcal{W}_3},\;
I_{\mathcal{W}_3}^{\mathrm{def}}\bigr),
\end{equation}
where the expected modular state spaces and partition function are
\begin{align}
% label removed: eq:w3-state-spaces
H_{g,n}^{\mathcal{W}_3}
&\;:=\;
\int_{\Sigma_{g,n}} V_{\mathcal{W}_3}^q,
\\
% label removed: eq:w3-partition-function
Z_{\mathcal{W}_3}(\h)
&\;:=\;
\sum_{g \geq 0} \h^{2g-2}\,\chi_{\gr}\!\bigl(H_g^{\mathcal{W}_3}\bigr).
\end{align}
\end{definition}

\begin{remark}[Organizing the three atoms]
% label removed: rem:three-atoms-organization
The three examples serve complementary roles in the modular holography programme:
\begin{enumerate}[label=\textup{(\roman*)}]
\item The free $\beta\gamma$ system is the \emph{strict free test}: all higher
  $A_\infty$ operations vanish, genus dependence enters only through determinant lines
  and factorization homology, and the modular $R$-matrix is exact at tree level.
\item Kodaira--Spencer holography is the \emph{holographic open/closed test}: the
  explicit universal defect algebra $U_{\mathrm{KS}}$ provides a concrete arena for
  testing open-closed modular splitting, and the single-trace decomposition gives a
  graded handle on the partition function.
\item The $\mathcal{W}_3$ model is the \emph{first nonlinear higher-spin test}: the
  nonlinear OPE $W \times W \sim T + \Lambda$ produces genuinely interacting modular
  operations, the conjectural slab reduction is expected to yield nontrivial effective central charges, and the
  modular envelope conjecture (Conjecture~\ref{conj:modular-envelope}) faces its first
  real obstruction-theoretic challenge.
\end{enumerate}
\end{remark}

%% ---------------------------------------------------------------
%%  Witten diagrams, N=4 matching, deformed conifold
%% ---------------------------------------------------------------

\subsection{Witten diagrams in twisted holography}
% label removed: subsec:witten-diagrams-twisted

In the twisted holographic setting, Witten diagrams admit a completely
algebraic reformulation that connects directly to the bar differential.

\begin{definition}[Twisted holographic propagator; \ClaimStatusProvedElsewhere]
% label removed: def:twisted-propagator
In holomorphic Chern--Simons theory on $\C \times \R_{>0}$ with gauge
algebra~$\mathfrak{g}$, the bulk-to-bulk propagator is
\begin{equation}% label removed: eq:hcs-propagator
K(z,w) \;=\; \frac{1}{z - w},
\end{equation}
while the boundary-to-bulk propagator is the Poisson kernel
$P(z,t;w) = t/\bigl((z-w)^2 + t^2\bigr)$.  These are the only
ingredients needed for all perturbative computations \cite{CG18}.
\end{definition}

\begin{proposition}[Two-point and three-point Witten diagrams; \ClaimStatusProvedElsewhere]
% label removed: prop:witten-2pt-3pt
The twisted Witten diagrams give \textup{(}\!\cite{CG18}, \S10\textup{)}:
\begin{enumerate}[label=\textup{(\roman*)}]
\item \textbf{Two-point function.}  A single propagator exchange yields
  \begin{equation}% label removed: eq:witten-2pt
  \bigl\langle O_a(z_1)\, O_b(z_2) \bigr\rangle
  \;=\;
  \frac{\kappa_{ab}}{(z_1 - z_2)^{2\Delta}},
  \end{equation}
  where $\kappa_{ab}$ is the Killing form and $\Delta$ is the conformal
  weight.
\item \textbf{Three-point function / OPE.}  The bulk cubic vertex coupled
  via three propagators gives
  \begin{equation}% label removed: eq:witten-3pt
  \bigl\langle O_a(z_1)\, O_b(z_2)\, O_c(z_3) \bigr\rangle
  \;=\;
  \frac{f^{d}_{ab}\,\kappa_{dc}}
       {(z_1 - z_2)^{\Delta_{12}}(z_2 - z_3)^{\Delta_{23}}(z_1 - z_3)^{\Delta_{13}}},
  \end{equation}
  recovering the structure constants $f^c_{ab}$ of~$\mathfrak{g}$.
\item \textbf{Stress tensor.}  In Kodaira--Spencer theory, the gravitational
  cubic vertex produces an OPE of stress-tensor type, reproducing the
  Virasoro algebra from Witten diagrams.
\end{enumerate}
\end{proposition}

\begin{remark}[Witten diagrams as the bar differential]
% label removed: rem:witten-bar-differential
The key structural insight is that \emph{Witten diagrams are Feynman
diagrams of the bar differential}: the bar complex $\bar{B}(\A)$ organizes
the full set of perturbative Witten diagrams into a nilpotent
differential.  The relation $d_{\bar{B}}^2 = 0$ is precisely the
statement that Witten diagram recursion is consistent: higher-point
correlators built from lower-point ones satisfy all crossing and
factorization identities.  This connects the holographic perturbation
theory of~\cite{CG18}, \S10, directly to the algebraic engine of
Part~I (Theorem~A).
\end{remark}

\begin{proposition}[Mode algebra matching; \ClaimStatusProvedElsewhere]
% label removed: prop:mode-algebra-matching
\index{mode algebra!holographic matching}
\textup{(Costello--Gaiotto \cite{CG18}, \S11.)}
The OPE of single-trace operators in the boundary chiral algebra is
reproduced by the perturbative expansion of the bulk holomorphic
Chern--Simons theory:
\begin{equation}% label removed: eq:mode-algebra-matching
  [O_{a,m},\, O_{b,n}]_{\mathrm{boundary}}
  \;=\;
  [O_{a,m},\, O_{b,n}]_{\mathrm{bulk\;Witten\;diagrams}}
\end{equation}
for all generators $O_{a,m}$ (where $a$ is a flavour index and $m$ a
mode number).  In particular:
\begin{enumerate}[label=\textup{(\roman*)},nosep]
\item The current-algebra modes $[J^a_m, J^b_n] = f^{c}_{ab}\,J^c_{m+n}
  + k\,m\,\delta_{m+n}\,\kappa^{ab}$ are reproduced by
  tree-level Witten diagrams with a single cubic vertex;
\item stress-tensor modes are reproduced by gravitational Witten
  diagrams in Kodaira--Spencer theory;
\item higher-spin modes in the $\mathcal{W}_\infty$ tower are
  reproduced by higher-point diagrams with the appropriate bulk
  vertices.
\end{enumerate}
\end{proposition}

\begin{remark}[Mode algebra matching as bar differential matching]
% label removed: rem:mode-algebra-bar
\index{mode algebra!bar differential}
The mode algebra matching \eqref{eq:mode-algebra-matching} is the
statement that the bar differentials of the boundary and bulk theories
agree: if $d_{\bar{B}}^{\mathrm{bdry}}$ is the bar differential of the
boundary chiral algebra and $d_{\bar{B}}^{\mathrm{bulk}}$ is the bar
differential of the bulk (i.e.\ the sum of all Witten diagram
contributions), then the matching is the quasi-isomorphism
\[
  \bar{B}(\cA_{\mathrm{boundary}})
  \;\xrightarrow{\;\sim\;}
  \bar{B}(\cA_{\mathrm{bulk}}).
\]
This is holographic Koszul duality at the chain level: not just an
isomorphism of cohomologies, but a quasi-isomorphism of dg coalgebras.
\end{remark}

\subsection{Global symmetry matching: the \texorpdfstring{$N\!=\!4$}{N=4} chiral algebra}
% label removed: subsec:n4-symmetry-matching

The most detailed test of the twisted holographic correspondence is the
complete matching of global symmetries between bulk and boundary for the
large-$N$ $\N=4$ chiral algebra.

\begin{definition}[Large-$N$ $\N=4$ chiral algebra; \ClaimStatusProvedElsewhere]
% label removed: def:large-n-n4-chiral
Following~\cite{CG18}, \S2, the large-$N$ $\N=4$ chiral algebra is the
BRST cohomology of a gauged $\beta\gamma$ system: start with
$N$ copies of $(b,c)$ ghosts and adjoint $\beta\gamma$ pairs
$(\beta_i, \gamma^i)$ for $i = 1,2,3$, with $\mathfrak{gl}_N$ gauge
symmetry.  At large~$N$, the single-trace sector provides an infinite
family of generators organized by spin.
\end{definition}

\begin{proposition}[Open-string symmetry matching; \ClaimStatusProvedElsewhere]
% label removed: prop:open-string-matching
The open-string global symmetry algebra of the boundary theory matches
the algebra of boundary operators in the holomorphic Chern--Simons bulk
\textup{(}\!\cite{CG18}, \S6.7\textup{)}.  Specifically:
\begin{enumerate}[label=\textup{(\roman*)}]
\item Every open-string symmetry generator in the boundary $\N=4$
  chiral algebra lifts to a bulk gauge transformation of the
  holomorphic Chern--Simons field;
\item the OPEs of the boundary generators are reproduced by
  bulk Witten diagrams \textup{(}Proposition~\textup{\ref{prop:witten-2pt-3pt})}.
\end{enumerate}
\end{proposition}

\begin{proposition}[Closed-string symmetry matching; \ClaimStatusProvedElsewhere]
% label removed: prop:closed-string-matching
The closed-string global symmetries match in both the bosonic and
fermionic sectors \textup{(}\!\cite{CG18}, \S\S6.8--6.10\textup{)}.
The bosonic closed-string symmetry algebra is identified with the
algebra of Kodaira--Spencer field configurations, while the fermionic
sector matches the ghost-number-shifted component of the single-trace
cohomology.
\end{proposition}

\begin{remark}[Affine quiver generalization]
% label removed: rem:affine-quiver-generalization
The construction extends to ADE affine quiver gauge theories
\textup{(}\!\cite{CG18}, \S2.4\textup{)}: for each ADE Dynkin
diagram~$\Gamma$, the corresponding quiver $\N=4$ theory produces a
chiral algebra $\A_\Gamma$ with the same holographic structure.  The
single-trace generators are organized by the representation theory
of~$\Gamma$, and the symmetry matching holds uniformly across the
ADE family.
\end{remark}

\begin{remark}[Symmetry matching as Koszul duality]
% label removed: rem:symmetry-koszul
In the language of Part~I, the global symmetry matching is the
statement
\begin{equation}% label removed: eq:symmetry-koszul-matching
\A_{\mathrm{bulk}} \;\simeq\; \A_{\partial}^{!},
\end{equation}
i.e., the bulk algebra is Koszul dual to the boundary algebra.  Each
boundary symmetry generator becomes a dual generator in the bulk
through Koszul duality: the bar differential encodes the
boundary OPEs, and the Verdier dual $D_{\mathrm{Ran}}(B(\cA))
\simeq B(\cA^!)$ produces the bar complex of the dual
boundary algebra~$\cA^!$, instantiating holographic Koszul duality
\textup{(}Conjecture~\textup{\ref*{conj:holographic-koszul-duality})}.
\end{remark}

\begin{theorem}[The Koszul triangle for $\mathcal{N}=4$;
\ClaimStatusConjectured]
% label removed: thm:n4-koszul-triangle
\index{Koszul triangle!N=4@$\mathcal{N}=4$}
For the large-$N$ $\mathcal{N}=4$ chiral algebra, the
bulk--boundary--line Koszul triangle
(Chapter~\ref{chap:ht-bulk-boundary-line}) specialises to
\begin{equation}% label removed: eq:n4-koszul-triangle
  \cA_{\mathrm{bulk}}
  \;\simeq\;
  Z_{\mathrm{der}}(\cA_{\partial})
  \;\simeq\;
  \operatorname{HH}^\bullet_{\mathrm{ch}}(\cA^!_{\mathrm{line}}),
\end{equation}
where $\cA_{\mathrm{bulk}}$ is holomorphic Chern--Simons on the
deformed conifold, $\cA_\partial$ is the $\mathcal{N}=4$ boundary VOA
(Definition~\ref{def:large-n-n4-chiral}), $\cA^!_{\mathrm{line}}$ is
the open-colour Koszul dual governing the line sector,
$Z_{\mathrm{der}}$ is the derived center, and
$\operatorname{HH}^\bullet_{\mathrm{ch}}$ is chiral Hochschild
cohomology.  Each vertex of the triangle has an independent
candidate description:
\begin{enumerate}[label=\textup{(\roman*)},nosep]
\item the bulk algebra is computed by KK reduction on $S^3$
  (Proposition~\ref{prop:kk-tower-s3});
\item the derived center of the boundary is computed by the open-string
  matching (Proposition~\ref{prop:open-string-matching});
\item the chiral Hochschild cohomology of the open-colour Koszul dual
  $\cA^!_{\mathrm{line}}$ is computed by the bar-cobar engine of
  Part~I (Theorem~H).
\end{enumerate}
The conjecture is that these three routes to $\cA_{\mathrm{bulk}}$
agree, which would provide one of the strongest available tests of
holographic Koszul duality.  In the Chriss--Ginzburg framework
(Proposition~\ref*{prop:chriss-ginzburg-structure} of Vol~I),
the bulk twisting morphism and the line-side chiral Hochschild
package are direct faces of a single master MC element
$\Theta_{\cN=4} \in
\operatorname{MC}(\mathfrak{g}^{\mathrm{amb}}_{\cN=4})$:
the bulk algebra is $\Theta\big|_{\hbar=0}$ (the twisting
morphism), and the chiral Hochschild cohomology is the
arity-preserving sub-MC
(Remark~\ref*{rem:five-from-one}, row~H).  The derived-center
vertex is instead the comparison corner suggested by the mixed part
of the master equation rather than a bare additional projection.
Thus the expected agreement of the three routes is organized by the master
equation, but the derived-center identification still uses exactly
the conjectural bulk-boundary comparison asserted in
\eqref{eq:n4-koszul-triangle}.  The master equation
$D\Theta_{\cN=4} + \tfrac{1}{2}[\Theta_{\cN=4},\Theta_{\cN=4}] = 0$.
\end{theorem}

\subsection{The deformed conifold: geometric origin of twisted holography}
% label removed: subsec:deformed-conifold

The geometric origin of the twisted holographic correspondence lies in
the deformed conifold, where B-brane backreaction produces the
holographic dual geometry.

\begin{proposition}[B-brane backreaction; \ClaimStatusProvedElsewhere]
% label removed: prop:b-brane-backreaction
Following~\cite{CG18}, \S\S4--5: placing $N$ B-branes at the tip of the
resolved conifold $T^*S^3$ produces a backreaction that deforms the
geometry.  The deformed conifold $\{xy - zw = \mu\} \subset \C^4$ with
deformation parameter $\mu \sim N\,g_s$ is the geometric dual of the
large-$N$ chiral algebra.
\end{proposition}

\begin{proposition}[Kaluza--Klein tower on $S^3$; \ClaimStatusProvedElsewhere]
% label removed: prop:kk-tower-s3
Kaluza--Klein reduction of the bulk fields on the $S^3$ fiber
decomposes them into towers indexed by $S^3$ harmonics
\textup{(}\!\cite{CG18}, \S\S4--5\textup{)}.  Each KK mode corresponds
to a single-trace generator of the boundary chiral algebra, giving a
higher-spin tower with one field at each integer spin $s \geq 1$.
\end{proposition}

\begin{remark}[Reconstruction of geometry from algebra]
% label removed: rem:geometry-from-algebra
A feature of the correspondence is that the deformed conifold
geometry is \emph{reconstructed from the algebraic structure} of the
global symmetry algebra \textup{(}\!\cite{CG18}, \S6.6\textup{)}: the
base and fiber of the conifold are read off from the representation
theory of the chiral algebra's mode algebra.
\end{remark}

\begin{remark}[The deformed conifold and the bar spectral sequence]
% label removed: rem:conifold-bar-spectral
In our framework, the deformed conifold geometry encodes the bar
complex's genus-tower data.  The KK decomposition on $S^3$ corresponds
to the $E_1$ page of the bar spectral sequence: each KK level
contributes a summand to $E_1^{p,q}(\bar{B}(\A))$, and the
differentials $d_r$ encode interactions between KK modes at successively
higher orders.  The convergence of the spectral sequence to
$H^*(\bar{B}(\A))$ is the algebraic analogue of the geometric statement
that the full backreacted geometry is determined by resumming all
KK interactions.
\end{remark}

\subsection{The celestial holographic modular Koszul datum}
% label removed: subsec:celestial-hmkd

\begin{construction}[Celestial holographic datum: first
principles; \ClaimStatusProvedHere]
% label removed: constr:celestial-hmkd
\index{celestial holography!holographic datum|textbf}
\index{holographic modular Koszul datum!celestial}
The celestial holographic modular Koszul datum
$\mathcal{H}(\mathrm{cel})$ is assembled from three inputs:
the celestial OPE algebra, its bar complex, and the collinear
splitting function.

\emph{Step~1: The celestial chiral algebra.}
Let $T$ be a $4$d massless gauge theory with holomorphic twist
(e.g.\ self-dual Yang--Mills or $\mathcal N\ge 1$ SYM).
The Mellin transform of the $S$-matrix produces a chiral
algebra $\mathcal{A}_{\mathrm{cel}}$ on
$\mathbb{CP}^1_{\mathrm{cel}}$ (the celestial sphere),
generated by conformally soft modes
$\mathcal{O}^a_{\Delta,J}$ of conformal dimension
$\Delta\in 1+\mathbb{Z}$ and spin~$J$.  The OPE
\[
\mathcal{O}^a_{\Delta_1,J_1}(z)\,
\mathcal{O}^b_{\Delta_2,J_2}(w)
\;\sim\;
\sum_{c,\Delta_3}\frac{C^{ab}{}_{c;\Delta_3}}{(z-w)^{\Delta_1+\Delta_2-\Delta_3}}\,
\mathcal{O}^c_{\Delta_3,J_1+J_2}(w)
\]
is determined by the collinear splitting amplitudes of~$T$.

\emph{Step~2: The ordered bar complex and celestial Yangian.}
The celestial chiral algebra is genuinely $E_1$ (the
collinear OPE is not commutative for non-abelian gauge
groups).  The ordered bar complex
${\Barch}^{\mathrm{ord}}(\mathcal{A}_{\mathrm{cel}})$ has bar
cohomology whose linear dual gives the open-colour dual
\[
\mathcal{A}_{\mathrm{cel}}^!_{\mathrm{line}}
\;=\;
H^*\!\bigl({\Barch}^{\mathrm{ord}}(\mathcal{A}_{\mathrm{cel}})\bigr)^\vee.
\]
One expects, and for
$\mathcal{A}_{\mathrm{cel}}=\widehat{\mathfrak{gl}}_N$
(the soft gluon sector of $\mathcal{N}=4$ SYM) one knows
\cite{Cos13}, that this open-colour dual is the celestial
Yangian
\[
\mathcal{A}_{\mathrm{cel}}^!_{\mathrm{line}}
\;\cong\;
Y_\hbar(\mathfrak{g}),
\]
where $\mathfrak{g}$ is the gauge Lie algebra and
$\hbar$ is determined by the coupling.  For
$\mathcal{A}_{\mathrm{cel}}=\widehat{\mathfrak{gl}}_N$
(the soft gluon sector of $\mathcal{N}=4$ SYM), this is
$Y(\mathfrak{gl}_N)$, the Yangian of $\mathfrak{gl}_N$.

\emph{Step~3: The collinear splitting function as classical
$r$-matrix.}
The tree-level collinear splitting function
$\mathrm{Split}_{+}(z_1,z_2)$ is the leading term
of the $R$-matrix:
\[
r_{\mathrm{cel}}(z)
\;=\;
\frac{\mathrm{Split}_{+}(z)}{z}
\;=\;
\frac{\Omega_{\mathfrak{g}}}{z}+O(z^{-2}),
\]
where $\Omega_{\mathfrak{g}}=\sum_a t^a\otimes t^a$ is the
Casimir of~$\mathfrak{g}$.  The classical Yang--Baxter
equation on $r_{\mathrm{cel}}(z)$
(Corollary~\ref{cor:classical-ybe}) is equivalent to the
soft-gluon collinear factorisation at tree level.

\emph{Step~4: The curvature and beta function.}
The chiral Koszul curvature is
\[
\kappa(\mathcal{A}_{\mathrm{cel}})
\;=\;
\frac{1}{2}\,\beta_1(T)\,\Omega_{\mathfrak{g}},
\]
where $\beta_1(T)$ is the one-loop beta function
coefficient.  At $\mathcal{N}=4$: $\beta_1=0$,
$\kappa=0$, and the shadow connection is flat.  For
$\mathcal{N}<4$: $\kappa\neq 0$, the shadow connection
acquires curvature, and the modular lift is obstructed
at genus~$1$ by $\mathrm{Obs}_1\propto\beta_1$.

\emph{Step~5: Assembly of the holographic modular Koszul datum.}
The six-fold datum
$\mathcal{H}(\mathrm{cel})
=(\mathcal{A},\,\mathcal{A}^!_{\mathrm{line}},\,C,\,r(z),\,
\Theta_{\mathcal{A}},\,\nabla^{\mathrm{hol}})$
is:
\begin{enumerate}[label=\textup{(\roman*)}]
\item $\mathcal{A}=\mathcal{A}_{\mathrm{cel}}$: the celestial
  OPE algebra of conformally soft modes.
\item $\mathcal{A}^!_{\mathrm{line}}
  =H^*\!\bigl({\Barch}^{\mathrm{ord}}(\mathcal{A}_{\mathrm{cel}})\bigr)^\vee$:
  the open-colour dual from ordered bar cohomology, conjecturally
  identified with the celestial Yangian $Y_\hbar(\mathfrak{g})$
  \textup{(}proved for $\widehat{\mathfrak{gl}}_N$ in
  $\mathcal{N}=4$ SYM\textup{)}.
\item $C={\Barch}^{\mathrm{ord}}(\mathcal{A}_{\mathrm{cel}})$:
  the ordered bar coalgebra.
\item $r(z)=\mathrm{Split}_{+}(z)/z$: the collinear
  splitting function as classical $r$-matrix.
\item $\Theta_{\mathcal{A}}$: the universal MC element in
  $\mathrm{MC}(\mathfrak{g}^{\mathrm{mod}}_{\mathcal{A}})$
  (from the bar-intrinsic construction, Vol~I
  Theorem~\ref*{V1-thm:mc2-bar-intrinsic}).
\item $\nabla^{\mathrm{hol}}=d-\mathrm{Sh}(\Theta_{\mathcal{A}})$:
  the shadow connection.  Flat iff $\kappa=0$
  ($\mathcal N=4$).
\end{enumerate}
The collinear limit is the collision residue:
$r_{\mathrm{cel}}(z)
=\operatorname{Res}^{\mathrm{coll}}_{0,2}
(\Theta_{\mathrm{cel}})$.
For $\mathcal{N}=4$: celestial amplitudes at $n$ points
satisfy the KZ differential equation
$\nabla^{\mathrm{hol}}_n\,\mathcal{M}_n=0$
with $\nabla^{\mathrm{hol}}_n
=d-\sum_{i<j}r_{\mathrm{cel}}(z_{ij})\,d\log z_{ij}$.
\end{construction}

\begin{computation}[Celestial holographic datum for
$\mathcal N=4$ SYM; \ClaimStatusConjectured]
% label removed: comp:celestial-hmkd-n4
\index{celestial holography!holographic datum}
For $\mathcal N=4$ SYM with gauge group $SU(N)$:
$\kappa=0$ (conformal), so
$\nabla^{\mathrm{hol}}$ is expected to be flat and
$\Theta_{\mathrm{cel}}$ is expected to lift to all genera.  One
expects the celestial Yangian $Y(\mathfrak{gl}_N)$ to control the line
operators; the corresponding R-matrix
$R_{\mathrm{cel}}(z)=1+\Omega/z+\cdots$ is expected to reproduce the
Parke--Taylor denominator structure of MHV amplitudes.
The modular $R$-matrix
$R^{\mathrm{mod}}(z;\hbar)=\sum_{g\ge 0}\hbar^{2g}r_g(z)$
is expected to encode genus corrections from loop amplitudes.
\end{computation}
